\documentclass{article}
\usepackage[utf8]{inputenc}
\usepackage{geometry}
\usepackage{enumitem}
\usepackage{hyperref}
\usepackage{xcolor}
\usepackage{graphicx}
\usepackage{array}
\usepackage{multirow}
\usepackage{amsmath, amssymb}
\usepackage{chemfig}
\usepackage{mhchem}
\usepackage{tikz}
\usepackage{setspace}
\usepackage{soul}
\usepackage{mathtools}
\usepackage{booktabs}
\usepackage{chemformula}
\usepackage{siunitx}
\usetikzlibrary{positioning, calc}
\usepackage{longtable}
\geometry{margin=1in}
\setlength{\parskip}{0.5em}
\usepackage{cancel}


\begin{document}

\begin{center}
    \vspace*{-1cm}
    \begin{tabular}{p{12cm} r}
        & \textbf{Name:} \underline{\hspace{5cm}} \\
    \end{tabular}

    \vspace{0.2cm}
    {\LARGE \textbf{Week 13 Worksheet}}
\end{center}    

\noindent\begin{tabular}{|p{0.9\linewidth}|}
\hline
\textbf{(1)} Which of the following changes how much a solution’s freezing point is lowered?\\[0.2cm]
\begin{minipage}[t]{0.85\linewidth}
\begin{enumerate}[label=\alph*.]
\item How much solute is added
\item The type of solute (number of particles it makes)
\item A \& B
\item The color of the solute
\end{enumerate}
\end{minipage}\\[0.2cm]
\hline
\end{tabular}

\noindent\begin{tabular}{|p{0.9\linewidth}|}
\hline
\textbf{(2)} What type of reaction occurs when methane burns in oxygen to produce carbon dioxide and water?\\[0.2cm]
CH$_4$ + 2O$_2$ $\rightarrow$ CO$_2$ + 2H$_2$O\\[0.2cm]
\begin{minipage}[t]{0.85\linewidth}
\begin{enumerate}[label=\alph*.]
\item Combination
\item Decomposition
\item Combustion
\item Single displacement
\end{enumerate}
\end{minipage}\\[0.2cm]
\hline
\end{tabular}

\noindent\begin{tabular}{|p{0.9\linewidth}|}
\hline
\textbf{(3)} What type of reaction occurs when potassium iodide reacts with lead(II) nitrate to form lead(II) iodide and potassium nitrate?\\[0.2cm]
2KI + Pb(NO$_3$)$_2$ $\rightarrow$ 2KNO$_3$ + PbI$_2$\\[0.2cm]
\begin{minipage}[t]{0.85\linewidth}
\begin{enumerate}[label=\alph*.]
\item Exchange
\item Combination
\item Decomposition
\item Combustion
\end{enumerate}
\end{minipage}\\[0.2cm]
\hline
\end{tabular}

\noindent\begin{tabular}{|p{0.9\linewidth}|}
\hline
\textbf{(4)} Vapor pressure lowering is an example of what type of property?\\[0.2cm]
\begin{minipage}[t]{0.85\linewidth}
\begin{enumerate}[label=\alph*.]
\item Colligative property
\item Chemical property
\item Physical property unrelated to concentration
\item Intensive property
\end{enumerate}
\end{minipage}\\[0.2cm]
\hline
\end{tabular}
\newpage

\noindent\begin{tabular}{|p{0.9\linewidth}|}
\hline
\textbf{(5)} Identify the reaction type:\\[0.2cm]
2H$_2$ + O$_2$ $\rightarrow$ 2H$_2$O\\[0.2cm]
\begin{minipage}[t]{0.85\linewidth}
\begin{enumerate}[label=\alph*.]
\item Exchage
\item Combination
\item Combustion
\item Decomposition
\end{enumerate}
\end{minipage}\\[0.2cm]
\hline
\end{tabular}

\noindent\begin{tabular}{|p{0.9\linewidth}|}
\hline
\textbf{(6)} In a redox reaction, the oxidizing agent is the species that:\\[0.2cm]
\begin{minipage}[t]{0.85\linewidth}
\begin{enumerate}[label=\alph*.]
\item Loses electrons
\item Gains electrons
\item Is always neutral
\item Has the highest atomic mass
\end{enumerate}
\end{minipage}\\[0.2cm]
\hline
\end{tabular}

\noindent\begin{tabular}{|p{0.9\linewidth}|}
\hline
\textbf{(7)} What type of chemical reaction occurs between sodium phosphate and calcium nitrate to form calcium phosphate and sodium nitrate?\\[0.2cm]
2Na$_3$PO$_4$ + 3Ca(NO$_3$)$_2$ $\rightarrow$ Ca$_3$(PO$_4$)$_2$ + 6NaNO$_3$\\[0.2cm]
\begin{minipage}[t]{0.85\linewidth}
\begin{enumerate}[label=\alph*.]
\item Decomposition
\item Exchange
\item Combination
\item Combustion
\end{enumerate}
\end{minipage}\\[0.2cm]
\hline
\end{tabular}

\noindent\begin{tabular}{|p{0.9\linewidth}|}
\hline
\textbf{(8)} Which of the following would have the highest osmolarity in water?\\[0.2cm]
\begin{minipage}[t]{0.85\linewidth}
\begin{enumerate}[label=\alph*.]
\item 1~M NaCl
\item 1~M glucose
\item 1~M MgCl$_2$
\item 1~M urea
\end{enumerate}
\end{minipage}\\[0.2cm]
\hline
\end{tabular}
\newpage


\noindent\begin{tabular}{|p{0.9\linewidth}|}
\hline
\textbf{(9)} A substance that gains electrons during a reaction is:\\[0.2cm]
\begin{minipage}[t]{0.85\linewidth}
\begin{enumerate}[label=\alph*.]
\item Oxidized
\item Reduced
\item Neutralized
\item Dissociated
\end{enumerate}
\end{minipage}\\[0.2cm]
\hline
\end{tabular}

\noindent\begin{tabular}{|p{0.9\linewidth}|}
\hline
\textbf{(10)} If two solutions have the same osmotic pressure, they are said to be:\\[0.2cm]
\begin{minipage}[t]{0.85\linewidth}
\begin{enumerate}[label=\alph*.]
\item Hypertonic
\item Hypotonic
\item Isotonic
\item Colloidal
\end{enumerate}
\end{minipage}\\[0.2cm]
\hline
\end{tabular}

\noindent\begin{tabular}{|p{0.9\linewidth}|}
\hline
\textbf{(11)} What type of reaction takes place when sodium carbonate reacts with calcium chloride to form calcium carbonate and sodium chloride?\\[0.2cm]
Na$_2$CO$_3$ + CaCl$_2$ $\rightarrow$ CaCO$_3$ + 2NaCl\\[0.2cm]
\begin{minipage}[t]{0.85\linewidth}
\begin{enumerate}[label=\alph*.]
\item Combination
\item Exchange
\item Single replacement
\item Decomposition
\end{enumerate}
\end{minipage}\\[0.2cm]
\hline
\end{tabular}

\noindent\begin{tabular}{|p{0.9\linewidth}|}
\hline
\textbf{(12)} In a redox reaction, the reducing agent is the substance that:\\[0.2cm]
\begin{minipage}[t]{0.85\linewidth}
\begin{enumerate}[label=\alph*.]
\item Gains electrons
\item Loses electrons
\item Is neither oxidized nor reduced
\item Lowers its oxidation number
\end{enumerate}
\end{minipage}\\[0.2cm]
\hline
\end{tabular}
\newpage


\noindent\begin{tabular}{|p{0.9\linewidth}|}
\hline
\textbf{(13)} What type of chemical reaction does ethanol undergo in air?\\[0.2cm]
C$_2$H$_5$OH + 3O$_2$ $\rightarrow$ 2CO$_2$ + 3H$_2$O\\[0.2cm]
\begin{minipage}[t]{0.85\linewidth}
\begin{enumerate}[label=\alph*.]
\item Neutralization
\item Single replacement
\item Combustion
\item Precipitation
\end{enumerate}
\end{minipage}\\[0.2cm]
\hline
\end{tabular}

\noindent\begin{tabular}{|p{0.9\linewidth}|}
\hline
\textbf{(14)} When a solute is added to a solvent, what happens to the freezing point of the solution?\\[0.2cm]
\begin{minipage}[t]{0.85\linewidth}
\begin{enumerate}[label=\alph*.]
\item It increases
\item It decreases
\item It stays the same
\item It becomes zero
\end{enumerate}
\end{minipage}\\[0.2cm]
\hline
\end{tabular}

\noindent\begin{tabular}{|p{0.9\linewidth}|}
\hline
\textbf{(15)} In the reaction Zn + Cu$^{2+}$ $\rightarrow$ Zn$^{2+}$ + Cu, which species is oxidized?\\[0.2cm]
\begin{minipage}[t]{0.85\linewidth}
\begin{enumerate}[label=\alph*.]
\item Cu$^{2+}$
\item Cu
\item Zn
\item H$_2$O
\end{enumerate}
\end{minipage}\\[0.2cm]
\hline
\end{tabular}

\noindent\begin{tabular}{|p{0.9\linewidth}|}
\hline
\textbf{(16)} Which factor determines how much the freezing point of a solution is lowered?\\[0.2cm]
\begin{minipage}[t]{0.85\linewidth}
\begin{enumerate}[label=\alph*.]
\item The number of solute particles
\item The solute's color
\item The solvent's taste
\item The solute's shape
\end{enumerate}
\end{minipage}\\[0.2cm]
\hline
\end{tabular}
\newpage


\noindent\begin{tabular}{|p{0.9\linewidth}|}
\hline
\textbf{(17)} What type of reaction occurs when acetylene reacts with oxygen to produce CO$_2$ and H$_2$O?\\[0.2cm]
2C$_2$H$_2$ + 5O$_2$ $\rightarrow$ 4CO$_2$ + 2H$_2$O\\[0.2cm]
\begin{minipage}[t]{0.85\linewidth}
\begin{enumerate}[label=\alph*.]
\item Neutralization
\item Combustion
\item Combination
\item Decomposition
\end{enumerate}
\end{minipage}\\[0.2cm]
\hline
\end{tabular}

\noindent\begin{tabular}{|p{0.9\linewidth}|}
\hline
\textbf{(18)} Adding salt to water causes which of the following effects?\\[0.2cm]
\begin{minipage}[t]{0.85\linewidth}
\begin{enumerate}[label=\alph*.]
\item Raises freezing point and lowers boiling point
\item Lowers freezing point and raises boiling point
\item Raises both freezing and boiling points
\item Lowers both freezing and boiling points
\end{enumerate}
\end{minipage}\\[0.2cm]
\hline
\end{tabular}

\noindent\begin{tabular}{|p{0.9\linewidth}|}
\hline
\textbf{(19)} What happens to the vapor pressure of a solvent when a nonvolatile solute is added?\\[0.2cm]
\begin{minipage}[t]{0.85\linewidth}
\begin{enumerate}[label=\alph*.]
\item It increases
\item It decreases
\item It stays the same
\item It becomes equal to the boiling point
\end{enumerate}
\end{minipage}\\[0.2cm]
\hline
\end{tabular}

\noindent\begin{tabular}{|p{0.9\linewidth}|}
\hline
\textbf{(20)} Identify the type of reaction below:\\[0.2cm]
K$_2$SO$_4$ + Ba(NO$_3$)$_2$ $\rightarrow$ 2KNO$_3$ + BaSO$_4$\\[0.2cm]
\begin{minipage}[t]{0.85\linewidth}
\begin{enumerate}[label=\alph*.]
\item Exchange
\item Combination
\item Combustion
\item Decomposition
\end{enumerate}
\end{minipage}\\[0.2cm]
\hline
\end{tabular}
\newpage


\noindent\begin{tabular}{|p{0.9\linewidth}|}
\hline
\textbf{(21)} Which solution would have the lowest vapor pressure?\\[0.2cm]
\begin{minipage}[t]{0.85\linewidth}
\begin{enumerate}[label=\alph*.]
\item Pure water
\item 1~M glucose solution
\item 1~M NaCl solution
\item 0.5~M NaCl solution
\end{enumerate}
\end{minipage}\\[0.2cm]
\hline
\end{tabular}

\noindent\begin{tabular}{|p{0.9\linewidth}|}
\hline
\textbf{(22)} Which of the following elements always has an oxidation number of +1 in compounds?\\[0.2cm]
\begin{minipage}[t]{0.85\linewidth}
\begin{enumerate}[label=\alph*.]
\item Sodium
\item Oxygen
\item Fluorine
\item Sulfur
\end{enumerate}
\end{minipage}\\[0.2cm]
\hline
\end{tabular}

\noindent\begin{tabular}{|p{0.9\linewidth}|}
\hline
\textbf{(23)} Identify the reaction type:\\[0.2cm]
2H$_2$O$_2$ $\rightarrow$ 2H$_2$O + O$_2$\\[0.2cm]
\begin{minipage}[t]{0.85\linewidth}
\begin{enumerate}[label=\alph*.]
\item Combination
\item Combustion
\item Decomposition
\item Exchange
\end{enumerate}
\end{minipage}\\[0.2cm]
\hline
\end{tabular}


\noindent\begin{tabular}{|p{0.9\linewidth}|}
\hline
\textbf{(24)} When hydrochloric acid reacts with sodium hydroxide to produce sodium chloride and water, what kind of reaction is it?\\[0.2cm]
HCl + NaOH $\rightarrow$  NaCl + H$_2$O\\[0.2cm]
\begin{minipage}[t]{0.85\linewidth}
\begin{enumerate}[label=\alph*.]
\item Combustion
\item Exchange
\item Combination
\item Decomposition
\end{enumerate}
\end{minipage}\\[0.2cm]
\hline
\end{tabular}
\newpage


\noindent\begin{tabular}{|p{0.9\linewidth}|}
\hline
\textbf{(25)} A 1.0~m solution of NaCl has a higher boiling point than a 1.0~m solution of glucose because:\\[0.2cm]
\begin{minipage}[t]{0.85\linewidth}
\begin{enumerate}[label=\alph*.]
\item NaCl has a smaller molar mass
\item NaCl dissociates into more particles
\item Glucose has more hydrogen bonds
\item Glucose is more polar
\end{enumerate}
\end{minipage}\\[0.2cm]
\hline
\end{tabular}

\noindent\begin{tabular}{|p{0.9\linewidth}|}
\hline
\textbf{(26)} What term describes the elevation of a solvent's boiling point when solute is added?\\[0.2cm]
\begin{minipage}[t]{0.85\linewidth}
\begin{enumerate}[label=\alph*.]
\item Freezing point depression
\item Boiling point elevation
\item Osmotic pressure
\item Vapor pressure lowering
\end{enumerate}
\end{minipage}\\[0.2cm]
\hline
\end{tabular}

\noindent\begin{tabular}{|p{0.9\linewidth}|}
\hline
\textbf{(27)} What kind of reaction is the burning of glucose in oxygen to form carbon dioxide and water?\\[0.2cm]
C$_6$H$_{12}$O$_6$ + 6O$_2$ $\rightarrow$  6CO$_2$ + 6H$_2$O\\[0.2cm]
\begin{minipage}[t]{0.85\linewidth}
\begin{enumerate}[label=\alph*.]
\item Double displacement
\item Combustion
\item Decomposition
\item Combination
\end{enumerate}
\end{minipage}\\[0.2cm]
\hline
\end{tabular}

\noindent\begin{tabular}{|p{0.9\linewidth}|}
\hline
\textbf{(28)} Which element is being reduced in the reaction Fe$_2$O$_3$ + 3CO $\rightarrow$ 2Fe + 3CO$_2$?\\[0.2cm]
\begin{minipage}[t]{0.85\linewidth}
\begin{enumerate}[label=\alph*.]
\item Fe
\item O
\item C
\item CO$_2$
\end{enumerate}
\end{minipage}\\[0.2cm]
\hline
\end{tabular}
\newpage


\noindent\begin{tabular}{|p{0.9\linewidth}|}
\hline
\textbf{(29)} When ammonium chloride reacts with silver nitrate to form silver chloride and ammonium nitrate, the reaction is classified as:\\[0.2cm]
NH$_4$Cl + AgNO$_3$ $\rightarrow$ AgCl + NH$_4$NO$_3$\\[0.2cm]
\begin{minipage}[t]{0.85\linewidth}
\begin{enumerate}[label=\alph*.]
\item Combination
\item Exchange
\item Decomposition
\item Combustion
\end{enumerate}
\end{minipage}\\[0.2cm]
\hline
\end{tabular}

\noindent\begin{tabular}{|p{0.9\linewidth}|}
\hline
\textbf{(30)} The decrease in vapor pressure when solute is added is known as:\\[0.2cm]
\begin{minipage}[t]{0.85\linewidth}
\begin{enumerate}[label=\alph*.]
\item Vapor pressure elevation
\item Vapor pressure lowering
\item Freezing point depression
\item Boiling point elevation
\end{enumerate}
\end{minipage}\\[0.2cm]
\hline
\end{tabular}

\noindent\begin{tabular}{|p{0.9\linewidth}|}
\hline
\textbf{(31)} When more solute is added to a solution, the osmotic pressure will:\\[0.2cm]
\begin{minipage}[t]{0.85\linewidth}
\begin{enumerate}[label=\alph*.]
\item Decrease
\item Stay the same
\item Increase
\item Become zero
\end{enumerate}
\end{minipage}\\[0.2cm]
\hline
\end{tabular}

\noindent\begin{tabular}{|p{0.9\linewidth}|}
\hline
\textbf{(32)} Osmotic pressure depends on which of the following factors?\\[0.2cm]
\begin{minipage}[t]{0.85\linewidth}
\begin{enumerate}[label=\alph*.]
\item The number of solute particles in solution
\item The color of the solute
\item The shape of the solvent molecules
\item The size of the container
\end{enumerate}
\end{minipage}\\[0.2cm]
\hline
\end{tabular}
\newpage


\noindent\begin{tabular}{|p{0.9\linewidth}|}
\hline
\textbf{(33)} What is the oxidation number of chlorine in NaCl?\\[0.2cm]
\begin{minipage}[t]{0.85\linewidth}
\begin{enumerate}[label=\alph*.]
\item +1
\item -1
\item 0
\item +2
\end{enumerate}
\end{minipage}\\[0.2cm]
\hline
\end{tabular}

\noindent\begin{tabular}{|p{0.9\linewidth}|}
\hline
\textbf{(34)} What type of reaction occurs when silver nitrate reacts with sodium chloride to produce silver chloride and sodium nitrate?\\[0.2cm]
AgNO$_3$ + NaCl $\rightarrow$ AgCl + NaNO$_3$\\[0.2cm]
\begin{minipage}[t]{0.85\linewidth}
\begin{enumerate}[label=\alph*.]
\item Combination
\item Decomposition
\item Exchange
\item Combustion
\end{enumerate}
\end{minipage}\\[0.2cm]
\hline
\end{tabular}

\noindent\begin{tabular}{|p{0.9\linewidth}|}
\hline
\textbf{(35)} What happens to the boiling point of a solvent when a nonvolatile solute is added?\\[0.2cm]
\begin{minipage}[t]{0.85\linewidth}
\begin{enumerate}[label=\alph*.]
\item It decreases
\item It remains the same
\item It increases
\item It fluctuates randomly
\end{enumerate}
\end{minipage}\\[0.2cm]
\hline
\end{tabular}

\noindent\begin{tabular}{|p{0.9\linewidth}|}
\hline
\textbf{(36)} When sodium reacts with chlorine to produce sodium chloride, what type of reaction is this?\\[0.2cm]
2Na + Cl$_2$ $\rightarrow$ 2NaCl\\[0.2cm]
\begin{minipage}[t]{0.85\linewidth}
\begin{enumerate}[label=\alph*.]
\item Decomposition
\item Combustion
\item Exchange
\item Combination
\end{enumerate}
\end{minipage}\\[0.2cm]
\hline
\end{tabular}
\newpage

\noindent\begin{tabular}{|p{0.9\linewidth}|}
\hline
\textbf{(37)} What happens to an atom's oxidation number when it is oxidized?\\[0.2cm]
\begin{minipage}[t]{0.85\linewidth}
\begin{enumerate}[label=\alph*.]
\item It increases
\item It decreases
\item It stays the same
\item It becomes zero
\end{enumerate}
\end{minipage}\\[0.2cm]
\hline
\end{tabular}

\noindent\begin{tabular}{|p{0.9\linewidth}|}
\hline
\textbf{(38)} What type of reaction occurs when lithium sulfate reacts with barium nitrate to form barium sulfate and lithium nitrate?\\[0.2cm]
Li$_2$SO$_4$ + Ba(NO$_3$)$_2$ $\rightarrow$ BaSO$_4$ + 2LiNO$_3$\\[0.2cm]
\begin{minipage}[t]{0.85\linewidth}
\begin{enumerate}[label=\alph*.]
\item Combination
\item Decomposition
\item Exchange
\item Combustion
\end{enumerate}
\end{minipage}\\[0.2cm]
\hline
\end{tabular}

\noindent\begin{tabular}{|p{0.9\linewidth}|}
\hline
\textbf{(39)} What is the oxidation number of nitrogen in NH$_3$?\\[0.2cm]
\begin{minipage}[t]{0.85\linewidth}
\begin{enumerate}[label=\alph*.]
\item +3
\item -3
\item 0
\item +5
\end{enumerate}
\end{minipage}\\[0.2cm]
\hline
\end{tabular}

\noindent\begin{tabular}{|p{0.9\linewidth}|}
\hline
\textbf{(40)} Identify the reaction type:\\[0.2cm]
BaCl$_2$ + Na$_2$SO$_4$ $\rightarrow$ BaSO$_4$ + 2NaCl\\[0.2cm]
\begin{minipage}[t]{0.85\linewidth}
\begin{enumerate}[label=\alph*.]
\item Combination
\item Decomposition
\item Exchange
\item Single replacement
\end{enumerate}
\end{minipage}\\[0.2cm]
\hline
\end{tabular}
\newpage


\noindent\begin{tabular}{|p{0.9\linewidth}|}
\hline
\textbf{(41)} What happens to osmotic pressure when temperature increases (all else equal)?\\[0.2cm]
\begin{minipage}[t]{0.85\linewidth}
\begin{enumerate}[label=\alph*.]
\item It decreases
\item It increases
\item It stays constant
\item It becomes zero
\end{enumerate}
\end{minipage}\\[0.2cm]
\hline
\end{tabular}

\noindent\begin{tabular}{|p{0.9\linewidth}|}
\hline
\textbf{(42)} The burning of benzene in air to form carbon dioxide and water is an example of what type of reaction?\\[0.2cm]
2C$_6$H$_6$ + 15O$_2$ $\rightarrow$ 12CO$_2$ + 6H$_2$O\\[0.2cm]
\begin{minipage}[t]{0.85\linewidth}
\begin{enumerate}[label=\alph*.]
\item Oxidation-reduction
\item Combustion
\item Decomposition
\item Neutralization
\end{enumerate}
\end{minipage}\\[0.2cm]
\hline
\end{tabular}

\noindent\begin{tabular}{|p{0.9\linewidth}|}
\hline
\textbf{(43)} The depression of the freezing point is directly proportional to which quantity?\\[0.2cm]
\begin{minipage}[t]{0.85\linewidth}
\begin{enumerate}[label=\alph*.]
\item Solute's molar mass
\item Solvent's density
\item Number of solute particles in solution
\item Volume of solvent
\end{enumerate}
\end{minipage}\\[0.2cm]
\hline
\end{tabular}

\noindent\begin{tabular}{|p{0.9\linewidth}|}
\hline
\textbf{(44)} Which factor mainly determines the amount of vapor pressure lowering in a solution?\\[0.2cm]
\begin{minipage}[t]{0.85\linewidth}
\begin{enumerate}[label=\alph*.]
\item The number of solute particles
\item The color of the solute
\item The shape of the solvent molecules
\item The pressure of the air above the solution
\end{enumerate}
\end{minipage}\\[0.2cm]
\hline
\end{tabular}
\newpage

\noindent\begin{tabular}{|p{0.9\linewidth}|}
\hline
\textbf{(45)} In the reaction 2Na + Cl$_2$ $\rightarrow$ 2NaCl, which is the oxidizing agent?\\[0.2cm]
\begin{minipage}[t]{0.85\linewidth}
\begin{enumerate}[label=\alph*.]
\item Na
\item Cl$_2$
\item NaCl
\item Both Na and Cl$_2$
\end{enumerate}
\end{minipage}\\[0.2cm]
\hline
\end{tabular}

\noindent\begin{tabular}{|p{0.9\linewidth}|}
\hline
\textbf{(46)} Osmosis is the movement of solvent molecules from:\\[0.2cm]
\begin{minipage}[t]{0.85\linewidth}
\begin{enumerate}[label=\alph*.]
\item Low solute concentration to high solute concentration
\item High solute concentration to low solute concentration
\item High pressure to low pressure
\item High temperature to low temperature
\end{enumerate}
\end{minipage}\\[0.2cm]
\hline
\end{tabular}

\noindent\begin{tabular}{|p{0.9\linewidth}|}
\hline
\textbf{(47)} What is osmotic pressure?\\[0.2cm]
\begin{minipage}[t]{0.85\linewidth}
\begin{enumerate}[label=\alph*.]
\item The pressure needed to stop osmosis
\item The pressure caused by boiling
\item The force between solute and solvent
\item The vapor pressure of a liquid
\end{enumerate}
\end{minipage}\\[0.2cm]
\hline
\end{tabular}

\noindent\begin{tabular}{|p{0.9\linewidth}|}
\hline
\textbf{(48)} Which of the following reactions is a redox reaction?\\[0.2cm]
\begin{enumerate}[label=\alph*.]
\item Zn + CuSO$_4$ $\rightarrow$ ZnSO$_4$ + Cu\\[0.1cm]
\item AgNO$_3$ + NaCl $\rightarrow$ AgCl + NaNO$_3$\\[0.1cm]
\item CaCO$_3$ $\rightarrow$ CaO + CO$_2$\\[0.1cm]
\item H$_2$O + SO$_3$ $\rightarrow$ H$_2$SO$_4$\\[0.1cm]
\end{enumerate}\\[0.2cm]
\hline
\end{tabular}
\newpage

\noindent\begin{tabular}{|p{0.9\linewidth}|}
\hline
\textbf{(49)} What happens to a cell placed in a hypotonic solution?\\[0.2cm]
\begin{minipage}[t]{0.85\linewidth}
\begin{enumerate}[label=\alph*.]
\item It shrinks
\item It swells
\item It stays the same size
\item It loses solute
\end{enumerate}
\end{minipage}\\[0.2cm]
\hline
\end{tabular}

\noindent\begin{tabular}{|p{0.9\linewidth}|}
\hline
\textbf{(50)} Which of the following increases osmotic pressure?\\[0.2cm]
\begin{minipage}[t]{0.85\linewidth}
\begin{enumerate}[label=\alph*.]
\item Lower solute concentration
\item Higher solute concentration
\item Larger solvent molecules
\item Higher freezing point
\end{enumerate}
\end{minipage}\\[0.2cm]
\hline
\end{tabular}

\noindent\begin{tabular}{|p{0.9\linewidth}|}
\hline
\textbf{(51)} Why does adding a solute lower the vapor pressure of a solvent?\\[0.2cm]
\begin{minipage}[t]{0.85\linewidth}
\begin{enumerate}[label=\alph*.]
\item Fewer solvent molecules can escape into the gas phase
\item Solute molecules turn into vapor more easily
\item The solute increases the temperature of the liquid
\item The solvent molecules become heavier
\end{enumerate}
\end{minipage}\\[0.2cm]
\hline
\end{tabular}

\noindent\begin{tabular}{|p{0.9\linewidth}|}
\hline
\textbf{(52)} What is the oxidation number of manganese in KMnO$_4$?\\[0.2cm]
\begin{minipage}[t]{0.85\linewidth}
\begin{enumerate}[label=\alph*.]
\item +2
\item +4
\item +6
\item +7
\end{enumerate}
\end{minipage}\\[0.2cm]
\hline
\end{tabular}
\newpage

\noindent\begin{tabular}{|p{0.9\linewidth}|}
\hline
\textbf{(53)} What is the oxidation number of sulfur in SO$_2$?\\[0.2cm]
\begin{minipage}[t]{0.85\linewidth}
\begin{enumerate}[label=\alph*.]
\item +2
\item +4
\item +6
\item 0
\end{enumerate}
\end{minipage}\\[0.2cm]
\hline
\end{tabular}

\noindent\begin{tabular}{|p{0.9\linewidth}|}
\hline
\textbf{(54)} If 0.5~mol of NaCl is dissolved in 1.0~kg of water, how many moles of particles are present assuming complete dissociation?\\[0.2cm]
\begin{minipage}[t]{0.85\linewidth}
\begin{enumerate}[label=\alph*.]
\item 0.5
\item 1.0
\item 1.5
\item 2.0
\end{enumerate}
\end{minipage}\\[0.2cm]
\hline
\end{tabular}

\noindent\begin{tabular}{|p{0.9\linewidth}|}
\hline
\textbf{(55)} The addition of a solute to a solvent causes the boiling point to:\\[0.2cm]
\begin{minipage}[t]{0.85\linewidth}
\begin{enumerate}[label=\alph*.]
\item Decrease
\item Increase
\item Remain constant
\item Become equal to the freezing point
\end{enumerate}
\end{minipage}\\[0.2cm]
\hline
\end{tabular}

\noindent\begin{tabular}{|p{0.9\linewidth}|}
\hline
\textbf{(56)} Which of the following represents a redox reaction?\\[0.2cm]
\begin{enumerate}[label=\alph*.]
\item HCl + NaOH $\rightarrow$ NaCl + H$_2$O\\[0.1cm]
\item Zn + 2HCl $\rightarrow$ ZnCl$_2$ + H$_2$\\[0.1cm]
\item CaCO$_3$ $\rightarrow$ CaO + CO$_2$\\[0.1cm]
\item NaOH + HNO$_3$ $\rightarrow$ NaNO$_3$ + H$_2$O\\[0.1cm]
\end{enumerate}\\[0.2cm]
\hline
\end{tabular}
\newpage

\noindent\begin{tabular}{|p{0.9\linewidth}|}
\hline
\textbf{(57)} Which of the following would cause the largest increase in the boiling point of water?\\[0.2cm]
\begin{minipage}[t]{0.85\linewidth}
\begin{enumerate}[label=\alph*.]
\item 1~mol NaCl
\item 1~mol glucose
\item 1~mol MgCl$_2$
\item 1~mol urea
\end{enumerate}
\end{minipage}\\[0.2cm]
\hline
\end{tabular}

\noindent\begin{tabular}{|p{0.9\linewidth}|}
\hline
\textbf{(58)} What is the oxidation number of oxygen in H$_2$O?\\[0.2cm]
\begin{minipage}[t]{0.85\linewidth}
\begin{enumerate}[label=\alph*.]
\item 0
\item -1
\item -2
\item +2
\end{enumerate}
\end{minipage}\\[0.2cm]
\hline
\end{tabular}

\noindent\begin{tabular}{|p{0.9\linewidth}|}
\hline
\textbf{(59)} The reaction between calcium hydroxide and hydrochloric acid produces calcium chloride and water. Identify the reaction type.\\[0.2cm]
Ca(OH)$_2$ + 2HCl $\rightarrow$ CaCl$_2$ + 2H$_2$O\\[0.2cm]
\begin{minipage}[t]{0.85\linewidth}
\begin{enumerate}[label=\alph*.]
\item Decomposition
\item Exchange
\item Combustion
\item Combination
\end{enumerate}
\end{minipage}\\[0.2cm]
\hline
\end{tabular}

\noindent\begin{tabular}{|p{0.9\linewidth}|}
\hline
\textbf{(60)} Osmolarity refers to:\\[0.2cm]
\begin{minipage}[t]{0.85\linewidth}
\begin{enumerate}[label=\alph*.]
\item The number of solute particles per liter of solution
\item The mass of solute per liter of solvent
\item The number of moles per kilogram of solvent
\item The ratio of solvent to solute molecules
\end{enumerate}
\end{minipage}\\[0.2cm]
\hline
\end{tabular}

\end{document}
